% GENERATED FILE. Edit canonical lessons, then run npm run generate:books.
% canonical-source-hash: fnv1a64:e4e68c53a76e3f6c
% canonical-lessons: KA-C519-udyami, KA-C519-vyavasthapaka, KA-C519-sahayaka, KA-C519-niyojisu, KA-C519-kudidu

\chapter{Running a Business}
\label{ch:ka-running-a-business}
\hlchaptermodality{519}

\begin{chapteropening}
\textbf{By the end of this chapter:} \emph{I can say an entrepreneur, a manager, an assistant, to assign and to save up.}

Complete the last lesson of chapter 519: I can say an entrepreneur, a manager, an assistant, to assign and to save up.
\end{chapteropening}

\section[\texorpdfstring{\kn{ಉದ್ಯಮಿ} (udyami)}{ಉದ್ಯಮಿ (udyami)}]{\kn{ಉದ್ಯಮಿ} (udyami) — an entrepreneur; an industrialist}

\label{lesson:KA-C519-udyami}



\begin{quote}
Before the new one: say the Kannada for retirement, then the Kannada for a contract.
\end{quote}

\subsection*{You'll want to know: \kn{ಉದ್ಯಮಿ}}
\textbf{\kn{ಉದ್ಯಮಿ}} — \emph{udyami} — \textquotedblleft{}an entrepreneur; an industrialist\textquotedblright{}.

\begin{quote}
\kn{ಆ} \kn{ಯುವ} \kn{ಉದ್ಯಮಿ} \kn{ಹೊಸ} \kn{ಕಾರ್ಖಾನೆ} \kn{ತೆರೆದನು}.
\end{quote}

\begin{quote}
\emph{ā yuva udyami hosa kārkhāne teredanu.}
\end{quote}

\begin{quote}
That young entrepreneur opened a new factory.
\end{quote}

Plural: \kn{ಉದ್ಯಮಿಗಳು}, \emph{udyamigaḷu}.

\kn{ಉದ್ಯಮ}, an industry or enterprise.

\begin{cousinweb}
\textbf{\kn{ಉದ್ಯಮಿ}} (\emph{udyami}) is borrowed from Sanskrit \emph{udyamin}, \textquotedblleft{}undertaking, persevering\textquotedblright{} (Monier-Williams), built on \emph{udyama}, \textquotedblleft{}undertaking, striving after\textquotedblright{} (CDIAL 2048).
\end{cousinweb}

\begin{grammarlens}[title={What you've built}]
One more word for describing what you lived and felt.
\end{grammarlens}

\subsection*{Your turn}
\begin{itemize}
\raggedright
  \item \emph{Say it:} \emph{udyami}
  \item \emph{Say it:} \emph{udyami}, once more
  \item \emph{Say it:} \emph{ā yuva udyami hosa kārkhāne teredanu.} — That young entrepreneur opened a new factory.
  \item \emph{Recall:} say the Kannada for retirement, then the Kannada for a contract, then say \emph{udyami} again
\end{itemize}

\begin{tcolorbox}[breakable,colback=teal!4,colframe=teal!35!black,title={Before you move on}]
What does \kn{ಉದ್ಯಮಿ} mean? (\textquotedblleft{}An entrepreneur; an industrialist\textquotedblright{}.) Say it once more.
\end{tcolorbox}
\section[\texorpdfstring{\kn{ವ್ಯವಸ್ಥಾಪಕ} (vyavasthāpaka)}{ವ್ಯವಸ್ಥಾಪಕ (vyavasthāpaka)}]{\kn{ವ್ಯವಸ್ಥಾಪಕ} (vyavasthāpaka) — a manager}

\label{lesson:KA-C519-vyavasthapaka}



\begin{quote}
Before the new one: say the Kannada for a contract, then the Kannada for an entrepreneur.
\end{quote}

\subsection*{You'll want to know: \kn{ವ್ಯವಸ್ಥಾಪಕ}}
\textbf{\kn{ವ್ಯವಸ್ಥಾಪಕ}} — \emph{vyavasthāpaka} — \textquotedblleft{}a manager\textquotedblright{}.

\begin{quote}
\kn{ಬ್ಯಾಂಕಿನ} \kn{ವ್ಯವಸ್ಥಾಪಕರು} \kn{ನನ್ನ} \kn{ಅರ್ಜಿಯನ್ನು} \kn{ಪರಿಶೀಲಿಸಿದರು}.
\end{quote}

\begin{quote}
\emph{byāṅkina vyavasthāpakaru nanna arjiyannu pariśīlisidaru.}
\end{quote}

\begin{quote}
The bank manager checked my application.
\end{quote}

Everyday talk also says \kn{ಮ್ಯಾನೇಜರ್}. \kn{ವ್ಯವಸ್ಥೆ} is an arrangement.

\begin{cousinweb}
\textbf{\kn{ವ್ಯವಸ್ಥಾಪಕ}} (\emph{vyavasthāpaka}) is borrowed from Sanskrit \emph{vyavasthāpaka}, \textquotedblleft{}settling, arranging, establishing\textquotedblright{} (Monier-Williams): the one who puts things in order.
\end{cousinweb}

\begin{grammarlens}[title={What you've built}]
One more word for describing what you lived and felt.
\end{grammarlens}

\subsection*{Your turn}
\begin{itemize}
\raggedright
  \item \emph{Say it:} \emph{vyavasthāpaka}
  \item \emph{Say it:} \emph{vyavasthāpaka}, once more
  \item \emph{Say it:} \emph{byāṅkina vyavasthāpakaru nanna arjiyannu pariśīlisidaru.} — The bank manager checked my application.
  \item \emph{Recall:} say the Kannada for a contract, then the Kannada for an entrepreneur, then say \emph{vyavasthāpaka} again
\end{itemize}

\begin{tcolorbox}[breakable,colback=teal!4,colframe=teal!35!black,title={Before you move on}]
What does \kn{ವ್ಯವಸ್ಥಾಪಕ} mean? (\textquotedblleft{}A manager\textquotedblright{}.) Say it once more.
\end{tcolorbox}
\section[\texorpdfstring{\kn{ಸಹಾಯಕ} (sahāyaka)}{ಸಹಾಯಕ (sahāyaka)}]{\kn{ಸಹಾಯಕ} (sahāyaka) — an assistant; a helper}

\label{lesson:KA-C519-sahayaka}



\begin{quote}
Before the new one: say the Kannada for an entrepreneur, then the Kannada for a manager.
\end{quote}

\subsection*{You'll want to know: \kn{ಸಹಾಯಕ}}
\textbf{\kn{ಸಹಾಯಕ}} — \emph{sahāyaka} — \textquotedblleft{}an assistant; a helper\textquotedblright{}.

\begin{quote}
\kn{ವೈದ್ಯರ} \kn{ಸಹಾಯಕ} \kn{ನನ್ನ} \kn{ಹೆಸರು} \kn{ಕೇಳಿದನು}.
\end{quote}

\begin{quote}
\emph{vaidyara sahāyaka nanna hesaru kēḷidanu.}
\end{quote}

\begin{quote}
The doctor's assistant asked my name.
\end{quote}

Plural: \kn{ಸಹಾಯಕರು}, \emph{sahāyakaru}.

Built on \kn{ಸಹಾಯ}, help, as in \kn{ಸಹಾಯ} \kn{ಮಾಡು}.

\begin{cousinweb}
\textbf{\kn{ಸಹಾಯಕ}} (\emph{sahāyaka}) is borrowed from Sanskrit \emph{sahāyaka}, the same as \emph{sahāya}, \textquotedblleft{}companion, helper\textquotedblright{} (CDIAL 13309).
\end{cousinweb}

\begin{grammarlens}[title={What you've built}]
One more word for describing what you lived and felt.
\end{grammarlens}

\subsection*{Your turn}
\begin{itemize}
\raggedright
  \item \emph{Say it:} \emph{sahāyaka}
  \item \emph{Say it:} \emph{sahāyaka}, once more
  \item \emph{Say it:} \emph{vaidyara sahāyaka nanna hesaru kēḷidanu.} — The doctor's assistant asked my name.
  \item \emph{Recall:} say the Kannada for an entrepreneur, then the Kannada for a manager, then say \emph{sahāyaka} again
\end{itemize}

\begin{tcolorbox}[breakable,colback=teal!4,colframe=teal!35!black,title={Before you move on}]
What does \kn{ಸಹಾಯಕ} mean? (\textquotedblleft{}An assistant; a helper\textquotedblright{}.) Say it once more.
\end{tcolorbox}
\section[\texorpdfstring{\kn{ನಿಯೋಜಿಸು} (niyōjisu)}{ನಿಯೋಜಿಸು (niyōjisu)}]{\kn{ನಿಯೋಜಿಸು} (niyōjisu) — to assign; to post}

\label{lesson:KA-C519-niyojisu}



\begin{quote}
Before the new one: say the Kannada for a manager, then the Kannada for an assistant.
\end{quote}

\subsection*{You'll want to know: \kn{ನಿಯೋಜಿಸು}}
\textbf{\kn{ನಿಯೋಜಿಸು}} — \emph{niyōjisu} — \textquotedblleft{}to assign; to post\textquotedblright{}.

\begin{quote}
\kn{ಕಂಪನಿ} \kn{ಅವಳನ್ನು} \kn{ಹೊಸ} \kn{ಶಾಖೆಗೆ} \kn{ನಿಯೋಜಿಸಿತು}.
\end{quote}

\begin{quote}
\emph{kampani avaḷannu hosa śākhege niyōjisitu.}
\end{quote}

\begin{quote}
The company posted her to the new branch.
\end{quote}

Past (I …): \kn{ನಿಯೋಜಿಸಿದೆ}, \emph{niyōjiside}.

The person takes -\kn{ಅನ್ನು}, the post or task -\kn{ಗೆ}.

\begin{cousinweb}
\textbf{\kn{ನಿಯೋಜಿಸು}} (\emph{niyōjisu}) is built on Sanskrit \emph{niyojana}, \textquotedblleft{}enjoining, directing, appointing to\textquotedblright{} (Monier-Williams), with the verb ending -\emph{isu}.
\end{cousinweb}

\begin{grammarlens}[title={What you've built}]
One more word for describing what you lived and felt.
\end{grammarlens}

\subsection*{Your turn}
\begin{itemize}
\raggedright
  \item \emph{Say it:} \emph{niyōjisu}
  \item \emph{Say it:} \emph{niyōjisu}, once more
  \item \emph{Say it:} \emph{kampani avaḷannu hosa śākhege niyōjisitu.} — The company posted her to the new branch.
  \item \emph{Recall:} say the Kannada for a manager, then the Kannada for an assistant, then say \emph{niyōjisu} again
\end{itemize}

\begin{tcolorbox}[breakable,colback=teal!4,colframe=teal!35!black,title={Before you move on}]
What does \kn{ನಿಯೋಜಿಸು} mean? (\textquotedblleft{}To assign; to post\textquotedblright{}.) Say it once more.
\end{tcolorbox}
\section[\texorpdfstring{\kn{ಕೂಡಿಡು} (kūḍiḍu)}{ಕೂಡಿಡು (kūḍiḍu)}]{\kn{ಕೂಡಿಡು} (kūḍiḍu) — to save up; to put by}

\label{lesson:KA-C519-kudidu}



\begin{quote}
Before the new one: say the Kannada for an assistant, then the Kannada for to assign.
\end{quote}

\subsection*{You'll want to know: \kn{ಕೂಡಿಡು}}
\textbf{\kn{ಕೂಡಿಡು}} — \emph{kūḍiḍu} — \textquotedblleft{}to save up; to put by\textquotedblright{}.

\begin{quote}
\kn{ಹೊಸ} \kn{ಚೀಲಕ್ಕಾಗಿ} \kn{ತಂಗಿ} \kn{ದುಡ್ಡು} \kn{ಕೂಡಿಟ್ಟಳು}.
\end{quote}

\begin{quote}
\emph{hosa cīlakkāgi taṅgi duḍḍu kūḍiṭṭaḷu.}
\end{quote}

\begin{quote}
My younger sister saved up money for a new bag.
\end{quote}

Past (I …): \kn{ಕೂಡಿಟ್ಟೆ}, \emph{kūḍiṭṭe}.

\kn{ಕೂಡು}, to gather + \kn{ಇಡು}, to put: to gather together and keep.

\begin{cousinweb}
\textbf{\kn{ಕೂಡಿಡು}} (\emph{kūḍiḍu}) is a native Dravidian compound: \emph{kūḍu}, \textquotedblleft{}to join, come together\textquotedblright{} (DEDR 1882), and \emph{iḍu}, \textquotedblleft{}to put, keep\textquotedblright{} (DEDR 442).
\end{cousinweb}

\begin{grammarlens}[title={What you've built}]
One more word for describing what you lived and felt.
\end{grammarlens}

\subsection*{Your turn}
\begin{itemize}
\raggedright
  \item \emph{Say it:} \emph{kūḍiḍu}
  \item \emph{Say it:} \emph{kūḍiḍu}, once more
  \item \emph{Say it:} \emph{hosa cīlakkāgi taṅgi duḍḍu kūḍiṭṭaḷu.} — My younger sister saved up money for a new bag.
  \item \emph{Recall:} say the Kannada for an assistant, then the Kannada for to assign, then say \emph{kūḍiḍu} again
\end{itemize}

\begin{tcolorbox}[breakable,colback=teal!4,colframe=teal!35!black,title={Before you move on}]
What does \kn{ಕೂಡಿಡು} mean? (\textquotedblleft{}To save up; to put by\textquotedblright{}.) Say it once more.
\end{tcolorbox}
